% !TeX program = xelatex
% ENGINE: xelatex (KHONG phai pdflatex). Ly do: refs.bib chua nguyen van ten
% van ban phap luat Viet Nam (QD 4801/QD-BNNMT, QD 1490/QD-TTg) nen phai render
% duoc dau tieng Viet. Da do: TeX Gyre Termes KHONG co dai U+1E00 (fc-list
% :lang=vi tra NO), con Charis SIL co (tra YES). Doi engine de giu trung thuc
% ten van ban, khong de escape tay \^{e}\`{} de sai.
% Paper 2 (IEEE 2-column) — style inherited from the senior co-author's
% "Toward Zero-Touch Smart Agriculture" (measured: 6527 words, 40 numbered
% equations, 8 tables, 0 figures, technical sections 148-630 words because the
% equations carry the content).
% Numbers: paper2/tables/*.tex are GENERATED from outputs/fraud_raw.csv by
% experiments/make_tables_paper2.py. Do not retype any number by hand.
\documentclass[journal]{IEEEtran}

% CHAP NHAN 4 canh bao 'Font shape TU/ptm/{m,bx}/{n,it} undefined'. Chung
% phat sinh TU BEN TRONG IEEEtran.cls (dong 503 va 1090), khi class tu chay
% \normalfont\selectfont voi rmdefault=ptm, tuc la TRUOC khi fontspec duoc
% nap -> moi sua o preamble deu la code chet. Da thu va DO LUONG hai cach:
%   (a) DeclareFontShape trong preamble   -> canh bao van con 4/4.
%   (b) DeclareFontShape TRUOC documentclass -> 4 loi FATAL 'Font family
%       TU+ptm unknown' (khai bao shape cho family chua ton tai). Da revert.
% Chung vo hai: luc do chua co chu nao duoc xuat, va pdffonts chung minh
% than bai chi nhung CharisSIL (4 bien the) + CM math; LMMono8 chi cua \url
% theo dung chuan IEEE. Xac minh font bang pdffonts, KHONG bang so canh bao
% trong log.

\usepackage{fontspec}
% Charis SIL: serif ho Palatino, du dau tieng Viet (SIL OFL). Fallback neu thieu.
\IfFontExistsTF{Charis SIL}{%
  \setmainfont{Charis SIL}%
}{%
  \IfFontExistsTF{Linux Libertine O}{%
    \setmainfont{Linux Libertine O}%
  }{%
    \setmainfont{DejaVu Serif}%
  }%
}

\usepackage{amsmath,amssymb,amsfonts}
\usepackage{amsthm}
\usepackage{booktabs}
\usepackage{graphicx}
\usepackage[hyphens]{url}
\usepackage[hidelinks]{hyperref}
\usepackage{cite}

\newtheorem{proposition}{Proposition}
\newtheorem{corollary}{Corollary}
\newtheorem{remark}{Remark}
\renewcommand{\proofname}{Proof}

% tight section spacing, as in the reference paper
\usepackage{titlesec}
\titlespacing{\section}{0pt}{6pt plus 2pt minus 2pt}{3pt plus 1pt minus 1pt}
\titlespacing{\subsection}{0pt}{5pt plus 2pt minus 2pt}{2pt plus 1pt minus 1pt}

% Kho cot that cua IEEEtran journal, do bang probe: 252pt = 8.85cm cho mot
% cot va 516pt = 17.8cm cho table*. tabcolsep mac dinh 6pt an het ngan sach
% cua bang 6-7 cot, nen dat 3pt de cac spec p{...} vua khuon.
\setlength{\tabcolsep}{3pt}

\begin{document}

\title{Physics-Attested Irrigation Logs: A Near-Zero-Cost MRV Layer for
Low-Emission Rice Under Alternate Wetting and Drying}

\author{Xuan~Van~Mai\IEEEauthorrefmark{1}\IEEEauthorrefmark{3}
        and~Tuong~Tri~Nguyen\IEEEauthorrefmark{2}\IEEEauthorrefmark{3}\IEEEauthorrefmark{4}%
\thanks{This work was supported by Hue University under Grant No. DHH2025-19-07.}%
\thanks{\IEEEauthorrefmark{1}Thuan Hoa High School, University of Education, Hue University, Hue, Vietnam.}%
\thanks{\IEEEauthorrefmark{2}Institute of Open Education and Information Technology, Hue University, Hue, Vietnam.}%
\thanks{\IEEEauthorrefmark{3}University of Education, Hue University, 34 Le Loi St., Hue, Vietnam.}%
\thanks{\IEEEauthorrefmark{4}Corresponding author: Tuong Tri Nguyen (e-mail: ntuongtri@hueuni.edu.vn).}}

\maketitle

\begin{abstract}
Measurement, reporting and verification (MRV) is the bottleneck keeping
emission reductions in low-emission rice from becoming tradable credits;
monitoring is about three quarters of that cost. Every credit-eligible
water-management action is already an \emph{actuation command}, and a record of
actual irrigation and drainage actions is the highest-ranked evidence under
Decision 4801/QD-BNNMT. We ask whether a farmer-written irrigation log can be
trusted with no sensor deployed. The log is replayed through the FAO-56 water
balance under public ERA5 forcing, turning each claimed depth into a trajectory
and exposing three violations no compliant log can produce: irrigation while the
root zone is full, a depth exceeding available storage, and depletion outside
its feasible range. We prove that an AWD-compliant log never triggers a filter,
and that a depth of at most $d^{*}\!\cdot\!\mathrm{span}+\varepsilon$ never
raises an alarm, so correct behaviour cannot be punished. Omission is the one
fraud class that physics cannot see; a single flow meter at the cooperative
pumping station closes it exactly. Across $240$ audited logs at three Mekong
Delta stations over two climate years the filters raise $0$ false alarms and
detect $60/60$ of each fraud class, with a mean pump gap of $29.0$\,mm. Tier~0
costs about $\$0$--$1$ per hectare per season, roughly $1\%$ of conservative
credit revenue. Expressed in depletion fraction, the scheme is immune to delta
subsidence of $1.1$--$2.5$\,cm per year, which would move an absolute
$15$\,cm threshold by up to $83\%$ in one crediting period.
\end{abstract}

\begin{IEEEkeywords}
Digital twin, MRV, carbon credit, alternate wetting and drying, rice,
water balance, fraud detection, networked control system.
\end{IEEEkeywords}

\section{Introduction}

The Mekong Delta programme on one million hectares of high-quality,
low-emission rice (Decision 1490/QD-TTg) passed its first-phase target, with
more than $354\,000$\,ha under the practice by the end of 2025, a measured
reduction of $3$--$4$\,t\,CO$_2$e per hectare and farm income up
$13.4\%$~\cite{mae2026progress}. The International Rice Research Institute
nevertheless identifies the next bottleneck as cost rather than
agronomy: without a functioning MRV system, emission reductions that are real
at field level cannot be formalised into a financial asset~\cite{irri2026mrv}.
Five obstacles stand between a smallholder and a credit.

\begin{enumerate}
\item \emph{Monitoring dominates cost.} A cross-study cost decomposition of
carbon-farming MRV attributes about $77\%$ of total cost to monitoring,
$13.6\%$ to reporting and $9.4\%$ to verification~\cite{jem2026mrv}. Any
architecture whose first move is to buy a sensor per plot attacks the largest
term with the largest capital outlay.
\item \emph{Revenue per hectare is small.} The revised rice methodology issues
about $4$--$6$ credits per hectare, down from the $12$--$15$ allowed under the
invalidated CDM methodology~\cite{rmi2025rice}, and rice methane credits
traded at $\$17$--$18$/t in reported Indian offtakes and $\$25$/t under the
Joint Crediting Mechanism in Indonesia~\cite{upstream2026price}. The whole MRV
budget is therefore of order tens of dollars per hectare per season.
\item \emph{The evidence that counts is a record of actions, not a
measurement.} Decision 4801/QD-BNNMT ranks ``the record of actual irrigation
and drainage'' above plans, schedules, field notes and
declarations~\cite{qd4801}. A sensor produces a state; a log produces an
action record, and it is the action record that the rule prefers.
\item \emph{Plots are small and instruments must survive a season.} Rural
connectivity is lossy, maintenance capacity is thin, and a device must run on
battery through a whole cropping season~\cite{irri2026mrv,delacruz2022}.
\item \emph{The measurement datum itself moves.} Groundwater extraction has
subsided the delta by about $18$\,cm over $25$ years, currently at
$1.1$\,cm/yr on average and above $2.5$\,cm/yr locally~\cite{minderhoud2017}.
An AWD threshold stated as an absolute depth below the surface therefore
drifts within a crediting period.
\end{enumerate}

The observation that drives this paper is that the scarce resource here is not
information about the field but \emph{trust in a record}. In our companion
work on command gating~\cite{twingate2026} the same digital twin was used to
decide whether an irrigation command may be sent over a lossy channel. Here the
twin changes role: it is not a controller but an \emph{auditor}. A claimed
irrigation depth is an input to a physical model, so a fabricated log produces a
trajectory that the soil cannot produce. Physics becomes the verification body,
and it is free.

The contributions are as follows.

\begin{enumerate}
\item \textbf{A zero-sensor attestation layer.} We formalise MRV evidence as an
irrigation log $L$ replayed through the FAO-56 root-zone balance under public
ERA5 forcing, and define three physical violation sets $V_1$, $V_2$, $V_4$ that
a compliant log cannot produce (Sec.~\ref{sec:filters}).
\item \textbf{Two guarantees that the detector cannot punish honest farmers.}
Proposition~\ref{prop:noalarm} shows that any log irrigating only at depletion
$d\ge d^{*}$ triggers no filter, and Proposition~\ref{prop:umax} gives the
closed-form largest claimable depth $u_{\max}=d^{*}\!\cdot\!\mathrm{span}
+\varepsilon$, evaluated at $61$\,mm for the reference soil
(Sec.~\ref{sec:guarantee}).
\item \textbf{An exact characterisation of what physics cannot see.} Omission
fraud leaves the log self-consistent, so no replay-based filter can detect it;
Proposition~\ref{prop:omission} shows that one flow meter at the cooperative
pumping station closes the gap exactly, and the measured gap is
$29.0$\,mm against a $5$\,mm decision threshold (Sec.~\ref{sec:omission}).
\item \textbf{Subsidence immunity by construction.} Because the balance is
written in depletion fraction relative to field capacity and wilting point, the
AWD criterion is invariant to a shifting surface elevation, whereas an absolute
$15$\,cm threshold drifts by $37$--$83\%$ of its own value within a five-year
crediting period (Sec.~\ref{sec:subsidence}).
\item \textbf{A tiered cost model with a break-even condition} and a
reproducible evaluation: $240$ audited logs, three stations, two climate years,
$0$ false alarms, $100\%$ detection of all three fraud classes
(Sec.~\ref{sec:cost}, Sec.~\ref{sec:results}).
\end{enumerate}

\section{System Model}
\label{sec:model}

\subsection{Field water balance}

The root-zone water content $w_t$ [mm] evolves hourly under the FAO-56 bucket
model~\cite{fao56}:
\begin{equation}
w_{t+1}=\mathrm{clip}_{\le FC}\!\left(w_t+P_t-RO_t-ET_{c,t}+u_t\right),
\label{eq:bucket}
\end{equation}
with crop evapotranspiration and a water-stress coefficient
\begin{equation}
ET_{c,t}=K_{s,t}\,K_c\,ET_{0,t},
\qquad
K_{s,t}=\frac{FC-w_t}{(1-p)\,\mathrm{TAW}},
\label{eq:etc}
\end{equation}
where $ET_{0,t}$ is the hourly reference evapotranspiration from the
Penman-Monteith equation, $P_t$ is rainfall, $RO_t$ is surface runoff from the
SCS curve-number method evaluated at daily resolution and shared pro rata over
the hours of a rain event, and $u_t$ is the irrigation depth applied in hour
$t$. Drainage above field capacity $FC$ is removed by the clipping operator, so
mass conservation holds exactly and is checked as an invariant. The root-zone
depletion fraction is
\begin{equation}
d_t=\frac{FC-w_t}{FC-WP},\qquad \mathrm{span}=FC-WP,
\label{eq:depl}
\end{equation}
with wilting point $WP$. For the reference silty clay loam of the delta,
$\mathrm{TAW}=70$\,mm, $FC=160$\,mm, $WP=90$\,mm, so $\mathrm{span}=70$\,mm.

\subsection{The irrigation log as evidence}

An irrigation log is a finite map
\begin{equation}
L=\{(t_i,u_i)\}_{i=1}^{|L|},\qquad u_i>0,
\label{eq:log}
\end{equation}
recorded by the farmer or by the gateway that actuates the pump. The
\emph{replay operator} $\mathcal{R}$ maps a log to the physical trajectory it
implies:
\begin{multline}
\mathcal{R}(L)=\{(\hat w_t,\hat d_t)\}_{t=0}^{T},\quad \hat w_0=w_{\mathrm{init}},\\
u_t=\sum_{(t_i,u_i)\in L}u_i\,\mathbb{1}[t_i{=}t],
\label{eq:replay}
\end{multline}
obtained by iterating \eqref{eq:bucket} with the observed forcing
$(P_t,ET_{0,t})$. Replay is deterministic and uses only public data, so any
auditor can reproduce it without access to the farm. A credit-eligible dry
phase is a maximal interval on which $\hat d_t\ge d_{\mathrm{dry}}$ and no
irrigation is claimed; Decision 4801 requires it to last at least three
days~\cite{qd4801}:
\begin{equation}
\begin{aligned}
\mathcal{A}(L)=\Big\{[a,b):\ &\min_{t\in[a,b)}\hat d_t\ge d_{\mathrm{dry}},\\
&u_t=0\ \forall t\in[a,b),\\
&b-a\ge 72\,\mathrm{h}\Big\}.
\end{aligned}
\label{eq:awd}
\end{equation}

\subsection{Threat model}

A fraudulent log is one that maximises $|\mathcal{A}(L)|$ or the claimed
management effort while the actual field history differs. Three classes cover
the ways a log can be altered.
\begin{enumerate}
\item \emph{Added events} ($F_{\mathrm{add}}$): claiming irrigation on hours
that were wet, to inflate the record of management activity or to disguise that
the field was continuously flooded.
\item \emph{Infeasible depth} ($F_{\mathrm{depth}}$): claiming a depth the root
zone cannot hold, to inflate reported water management or to match a subsidy
form.
\item \emph{Omission} ($F_{\mathrm{omit}}$): deleting real irrigation commands
inside a window so that the replayed trajectory shows an unbroken dry phase of
at least $72$\,h and a credit is claimed for drainage that did not happen. This
is the only class that is profitable \emph{and} invisible to replay, because a
log with entries removed remains physically self-consistent.
\end{enumerate}

\subsection{Attestation problem}

Given a log $L$, produce a decision $a(L)\in\{\text{accept},
\text{query},\text{reject}\}$ that minimises the cost of MRV subject to
two constraints: no compliant log is rejected, and every fraudulent log is
either queried or rejected. We impose the first constraint as a hard guarantee
(Prop.~\ref{prop:noalarm}) rather than as a measured rate, because a detector
that punishes honest farmers will not be adopted regardless of its accuracy.


\section{Physics-Based Log Filters}
\label{sec:filters}

A replayed log produces, at each claimed event $(t_i,u_i)$, the pre-event
state $\hat w_{t_i}^{-}$ and hence the storage the soil can still accept,
$S_i=FC-\hat w_{t_i}^{-}$. Three violation sets are read off the replay.

\emph{Saturation violation.} A claim made when the root zone is already at or
above field capacity cannot be stored:
\begin{equation}
V_1(L)=\Big\{(t_i,u_i)\in L:\ \hat w_{t_i}^{-}\ \ge\ FC-\varepsilon\Big\},
\label{eq:v1}
\end{equation}
where $\varepsilon>0$ is a tolerance against forcing error.

\emph{Infeasible depletion.} The replayed depletion must remain inside its
physical range; a claim history that drives the soil permanently below the
wilting point while irrigation is claimed is not realisable:
\begin{equation}
V_2(L)=\Big\{t:\ \hat d_t>1+\eta\Big\},\qquad \eta=0.02 .
\label{eq:v2}
\end{equation}
The tolerance $\eta$ is needed because \eqref{eq:bucket} integrates depletion
below $\mathrm{TAW}$ under sustained deficit rather than clamping it, which is
the FAO-56 convention; $\eta$ keeps the filter from firing on a legitimately
stressed field.

\emph{Storage excess.} The claimed depth is compared with the storage actually
available at that hour:
\begin{equation}
V_4(L)=\Big\{(t_i,u_i)\in L:\ u_i>S_i+\varepsilon\Big\}.
\label{eq:v4}
\end{equation}
$V_1$ is the limiting case of $V_4$ as $S_i\to0$; both are reported separately
so that a reviewer can see whether a rejection comes from a full profile or
from an oversized claim.

The detector is the union of the three,
\begin{equation}
\begin{aligned}
V(L)&=V_1(L)\cup V_2(L)\cup V_4(L),\\
a(L)&=
\begin{cases}
\text{accept}, & V(L)=\emptyset,\\
\text{query}, & 0<|V(L)|\le \kappa,\\
\text{reject}, & |V(L)|>\kappa.
\end{cases}
\end{aligned}
\label{eq:decision}
\end{equation}
with $\kappa$ set by the verifier's appetite for field visits. The decision is
a triage signal for the human auditor, not an automated penalty; the standard
still requires third-party verification~\cite{vm0051}.

\section{Guarantees Against False Accusation}
\label{sec:guarantee}

The adoption constraint of Sec.~\ref{sec:model} requires that a compliant log
never be rejected. We state it as a property of the filters rather than as an
empirical rate.

\begin{proposition}[Compliant logs are clean]
\label{prop:noalarm}
Let $L$ irrigate only at hours whose pre-event depletion satisfies
$\hat d_{t_i}^{-}\ge d^{*}$, with every claimed depth $u_i\le
d^{*}\!\cdot\!\mathrm{span}+\varepsilon$. Then $V_1(L)=V_2(L)=V_4(L)=\emptyset$.
\end{proposition}

\begin{proof}
$V_1$: $\hat d_{t_i}^{-}\ge d^{*}>0$ gives
$\hat w_{t_i}^{-}=FC-\hat d_{t_i}^{-}\,\mathrm{span}\le FC-d^{*}\mathrm{span}
<FC-\varepsilon$ for any $\varepsilon<d^{*}\mathrm{span}$, so \eqref{eq:v1}
cannot fire. $V_4$: the available storage is
$S_i=FC-\hat w_{t_i}^{-}=\hat d_{t_i}^{-}\,\mathrm{span}\ge
d^{*}\mathrm{span}$, hence $u_i\le d^{*}\mathrm{span}+\varepsilon\le
S_i+\varepsilon$ and \eqref{eq:v4} cannot fire. $V_2$: irrigation at
$\hat d\ge d^{*}$ with depth at most $d^{*}\mathrm{span}+\varepsilon$ raises
the water content by at most the deficit, so the post-event depletion is
non-negative; between events depletion rises only through $ET_c$, which is
bounded by the same forcing that produced the reference trajectory, hence
$\hat d_t\le 1+\eta$.
\end{proof}

The bound on the claimable depth is the quantity a farmer-facing application
must enforce, so we isolate it.

\begin{proposition}[Largest admissible claim]
\label{prop:umax}
For the reference profile and $d^{*}=0.80$, $\varepsilon=5$\,mm,
\begin{equation}
u_{\max}=d^{*}\!\cdot\!\mathrm{span}+\varepsilon=0.80\times70+5=61\ \text{mm},
\label{eq:umax}
\end{equation}
and a claim fires $V_4$ if and only if $u>61$\,mm at that depletion.
\end{proposition}

Table~\ref{tab:prop1} lists the derivation, and Sec.~\ref{sec:results} verifies
it by scanning claimed depths: $20$, $40$, $60$ and $61$\,mm pass while $62$,
$80$ and $160$\,mm fire, with no disagreement between the closed form and the
replay.

\begin{remark}
The agronomic candidate depths used in practice, $\{20,40,60\}$\,mm, all lie
below $u_{\max}$, so the filter never obstructs a normal AWD schedule. It fires
only on claims that exceed what a single irrigation event can physically
deliver to a root zone at the drainage threshold.
\end{remark}

\section{Omission and the Pump-Meter Closure}
\label{sec:omission}

\begin{proposition}[Replay is blind to omission]
\label{prop:omit-blind}
Let $L$ be a compliant log and $L'\subset L$ be obtained by deleting entries.
Then $V(L')=\emptyset$ whenever $V(L)=\emptyset$: removing claims cannot create
a physical violation.
\end{proposition}

\begin{proof}
Each filter in \eqref{eq:v1}--\eqref{eq:v4} is evaluated on a claimed event or
on a state reached under claimed events. Deleting an event removes terms from
$V_1$ and $V_4$ directly. For $V_2$, the replayed depletion with less claimed
water is pointwise no smaller than with more, and it is bounded above by the
reference trajectory that produced no violation, because the reference already
integrates the full deficit forcing. Hence no filter fires.
\end{proof}

This is not a defect to be hidden but the exact boundary of what a free,
physics-only layer can attest, and it determines the cheapest device that
closes it. Let $M$ be the total volume measured at the cooperative pumping
station over the season and $U(L)=\sum_{(t_i,u_i)\in L}u_i\,A$ the volume
implied by the log over the served area $A$. Define the pump gap
\begin{equation}
G=M-U(L).
\label{eq:gap}
\end{equation}

\begin{proposition}[One flow meter closes omission]
\label{prop:omission}
If a farmer omits events totalling $\Delta$ mm, then $G$ increases by exactly
$\Delta A$, independent of the weather, the soil and the timing of the omitted
events. Omission is therefore detected by the single scalar test
$G>\tau$ with $\tau$ a volume tolerance.
\end{proposition}

\begin{proof}
$M$ is measured, not claimed. $U(L)$ is linear in the claimed depths. Removing
$\Delta$ mm of claims reduces $U$ by $\Delta A$, so $G=M-U$ rises by exactly
$\Delta A$. No term in \eqref{eq:gap} depends on $P_t$, $ET_{0,t}$ or the
placement of the omitted events.
\end{proof}

Proposition~\ref{prop:omission} is the reason the tier structure of
Sec.~\ref{sec:cost} exists. A shared flow meter costs one device per
cooperative rather than one sensor per plot, and it is the unique minimal
addition that converts an attestable-but-incomplete record into a complete one.
In the evaluation the mean gap of omission frauds is $29.0$\,mm against a
$\tau=5$\,mm threshold, so the margin is nearly six-fold.

\section{Subsidence Invariance}
\label{sec:subsidence}

AWD compliance is usually stated as an absolute water-table or ponding depth,
for example $15$\,cm below the surface~\cite{carrijo2017}. In a subsiding delta
the surface elevation $z_s(t)$ falls at rate $\dot z$, so an absolute threshold
measured against a fixed benchmark drifts.

\begin{proposition}[Drift of an absolute threshold]
\label{prop:subsid}
Over a crediting period of $T_c$ years an absolute threshold expressed against
a fixed datum shifts by $\dot z\,T_c$. For $\dot z\in[1.1,2.5]$\,cm/yr and
$T_c=5$\,yr the shift is $5.5$--$12.5$\,cm, that is $37$--$83\%$ of a
$15$\,cm threshold. A criterion written in the depletion fraction
\eqref{eq:depl} has zero drift, because $FC$ and $WP$ are soil properties
referenced to the root zone and not to an absolute elevation.
\end{proposition}

\begin{proof}
The absolute criterion compares a measured depth to a constant benchmark; if
the benchmark and the surface move apart at $\dot z$, the effective threshold
changes by $\dot z T_c$ over $T_c$ years. The depletion criterion compares
$FC-w_t$ to $(FC-WP)$, both of which are defined within the root zone and
translate rigidly with the surface, so their ratio is invariant to $z_s$.
\end{proof}

The subsidence rates are those measured for the delta by
\cite{minderhoud2017}. The practical consequence is that a scheme anchored on
depletion fraction needs no re-surveying of field benchmarks during a five-year
crediting period, whereas an absolute-depth scheme does.


\section{Cost Model and Break-Even}
\label{sec:cost}

Let a cooperative serve $A$ hectares with $N$ cropping seasons of device life,
and let the MRV cost per hectare per season decompose as
\begin{equation}
C=\underbrace{\frac{n_{\mathrm{dev}}\,p_{\mathrm{dev}}}{A\,N}}_{\text{hardware}}
+\underbrace{c_{\mathrm{conn}}}_{\text{connectivity}}
+\underbrace{h\,w}_{\text{labour}}
+\underbrace{\frac{c_{\mathrm{audit}}}{A}}_{\text{verification}},
\label{eq:cost}
\end{equation}
where $n_{\mathrm{dev}}$ and $p_{\mathrm{dev}}$ are the number and unit price of
devices, $h$ is the labour hours per hectare per season and $w$ the rural wage.
Revenue per hectare is
\begin{equation}
R=\kappa\,\rho ,
\label{eq:revenue}
\end{equation}
with $\kappa$ credits issued per hectare and $\rho$ the price per credit. The
scheme is economically viable when
\begin{equation}
C\ \le\ \beta R ,\qquad \beta\in(0,1),
\label{eq:breakeven}
\end{equation}
and $\beta=0.5$ is the level at which monitoring consumes half of the credit
revenue. Substituting \eqref{eq:revenue} gives the break-even issuance
\begin{equation}
\kappa_{\min}=\frac{C}{\beta\rho}.
\label{eq:kappamin}
\end{equation}

The three tiers differ only in $n_{\mathrm{dev}}$ and in which fraud class they
close. Tier~0 sets $n_{\mathrm{dev}}=0$: the log is written on a handset the
farmer already owns, and the filters of Sec.~\ref{sec:filters} run on public
weather data at no marginal cost. Tier~1 adds a single flow meter at the
cooperative pumping station, so $n_{\mathrm{dev}}=1$ for the whole area $A$
rather than one device per plot; by Prop.~\ref{prop:omission} this one device is
sufficient for omission. Tier~2 adds one water-level node per cooperative,
which re-anchors the twin and restores the calibration property that a
measurement-based MRV would need anyway.

Table~\ref{tab:cost} reports the three tiers and Table~\ref{tab:baseline}
compares them with the three architectures they replace. With
$\kappa\in[4,6]$ and $\rho\in[\$15,\$25]$~\cite{rmi2025rice,upstream2026price},
revenue is $R\in[\$60,\$150]$ per hectare per season, so
$\kappa_{\min}$ at $\beta=0.5$ corresponds to a monitoring budget of
$\$30$--$75$ per hectare. The full three-tier stack sits an order of magnitude
below that ceiling, and tier~0 alone is about $1\%$ of conservative revenue.

\begin{remark}
The device prices in Table~\ref{tab:cost} are planning estimates, not
quotations. The ordering of the tiers and the break-even inequality
\eqref{eq:breakeven} do not depend on them; only the absolute margin does. A
quotation should replace the estimates before any procurement decision.
\end{remark}

\section{Experimental Design}
\label{sec:design}

\subsection{Forcing and reference trajectories}

Weather is hourly ERA5 reanalysis obtained from the public Open-Meteo archive
for three Mekong Delta stations, Can Tho, Soc Trang and Ca Mau, over two
climate years, 2024 (water-scarce) and 2025 (water-surplus), with no missing
values. The window is 1 March to 29 May, $2\,160$ hours, which is the dry
season in which AWD is practised. For each station, year and seed a reference
trajectory is generated by iterating \eqref{eq:bucket} under an AWD-compliant
schedule that irrigates whenever $\hat d\ge d^{*}=0.80$ with a depth drawn from
$\{20,40\}$\,mm. This reference plays the role that a sensor would play in a
deployed system: it is what actually happened.

\subsection{Fraud generation}

Three fraud classes are injected into the honest log $L$, each with a fixed
parameterisation so that the result is not tuned per class.

\emph{Added events.} Eight additional claims of $\{20,40,60\}$\,mm are inserted
at hours with rainfall above $5$\,mm that carry no honest claim, modelling a
farmer who inflates the record of management activity.

\emph{Infeasible depth.} One third of the honest entries are replaced by a claim
equal to the whole storage span, $u=FC$, which no single event can deliver.

\emph{Omission.} A $96$-hour window is selected, every honest claim inside it is
deleted, and the resulting log is presented as an unbroken dry phase, which is
the profitable fraud because it manufactures credit-eligible drainage.

\subsection{Decision rules under test}

Tier~0 applies $V(L)$ from \eqref{eq:decision}. Tier~1 applies the scalar test
$G>\tau$ with $\tau=5$\,mm, where $G$ is computed from
\eqref{eq:gap} using the reference total as the measured pump volume. Both are
evaluated on the same $240$ logs, with no filter and no threshold adjusted after
inspection of the outcomes.

\subsection{Metrics}

Two rates are reported per fraud class. The false-alarm rate on honest logs is
\begin{equation}
\mathrm{FAR}=\frac{1}{n_{\mathrm{h}}}\sum_{j=1}^{n_{\mathrm{h}}}\mathbb{1}\big[V(L_j)\neq\emptyset\big],
\label{eq:far}
\end{equation}
and the detection rate of class $c$ is
\begin{equation}
\mathrm{DR}_c=\frac{1}{n_c}\sum_{j=1}^{n_c}\mathbb{1}\big[\text{class }c\text{ flagged on }L_j\big].
\label{eq:dr}
\end{equation}
By the adoption constraint of Sec.~\ref{sec:model}, $\mathrm{FAR}$ must be
exactly zero; a nonzero value invalidates the scheme irrespective of
$\mathrm{DR}_c$. The experiment is scripted so that any nonzero $\mathrm{FAR}$
or any $\mathrm{DR}_c$ below $50\%$ returns a nonzero exit code.

\section{Numerical Results}
\label{sec:results}

Table~\ref{tab:setup} summarises the configuration. All $240$ logs were audited;
every number in this section is regenerated from the released raw output by the
same script that produced the tables.

\subsection{Detection performance}

Table~\ref{tab:detect} reports the outcome per station and year. The detector
raises \emph{no} false alarm on any of the $60$ honest logs and detects all
$60$ logs of each fraud class, in every one of the six station-year cells
individually rather than only in aggregate. Detection of added events is shared
between $V_1$ and $V_4$; detection of infeasible depth is carried by $V_4$
alone, since a $160$\,mm claim at depletion $0.80$ leaves the profile below
field capacity and therefore does not saturate it.

Omission, as predicted by Prop.~\ref{prop:omit-blind}, is invisible at tier~0
and is detected in $60/60$ cases at tier~1. This is the intended division of
labour rather than a failure of the physics layer.

\subsection{Intensity of the violations}

Table~\ref{tab:intensity} separates the two questions a verifier asks: how loud
is the signal, and how large is the margin. Added-event frauds trigger
$|V_1|=5.87$ and $|V_4|=7.25$ violations per log on average, and inflate the
claimed total from $124.7$\,mm to $445.0$\,mm, a factor of $3.6$. Infeasible
depth triggers $|V_4|=1.57$ per log because only a third of the entries were
altered. Omission leaves every physical filter silent, as required by
Prop.~\ref{prop:omit-blind}, but produces a positive pump gap of $29.0$\,mm
while honest logs produce $G=0.0$ exactly, against a $\tau=5$\,mm threshold.

The maximum depletion reached is also diagnostic. Honest logs peak at
$\bar d_{\max}=0.802$, consistent with irrigating at the $0.80$ threshold,
whereas omission frauds reach $0.896$ because the deleted irrigation never
arrives in the replay. A verifier can therefore use $\bar d_{\max}$ as a soft
indicator even where $G$ is unavailable, although we do not rely on it.

\subsection{Verification of the admissible-depth bound}

Prop.~\ref{prop:umax} predicts a hard boundary at $61$\,mm. We scanned claimed
depths through the replay operator at the drainage threshold: $20$, $40$, $60$
and $61$\,mm produced no $V_4$ violation, and $62$, $80$ and $160$\,mm each
produced one. The closed form and the simulation agree at every point, with no
disagreement to reconcile. Table~\ref{tab:prop1} lists the derivation.

\subsection{Summary of headline results}

Table~\ref{tab:headline} collects every headline number with a pointer to its
source, so that a reader can trace each claim to a table or a proposition
without searching the text.

\section{Relevance to MRV Practice}
\label{sec:practice}

The scheme is a layer that produces activity data, not a complete MRV system.
Table~\ref{tab:legal} maps each requirement of the two governing instruments to
the component that meets it, and states explicitly what is not claimed. Three
points deserve emphasis.

First, the evidence produced is of the highest rank in the ordering of Decision
4801, namely a record of actual irrigation and drainage actions, and it is at
hourly resolution, whereas the current procedure relies on periodic photographs.
Second, the uncertainty machinery is shared with the standard: the calibration
interval used in the companion work is the $95\%$ interval with $z=1.96$ that
Appendix A12 of Decision 4801 prescribes~\cite{qd4801}, and the conservative
bias of the thresholds in \eqref{eq:v1}--\eqref{eq:v4} points in the same
direction as the conservativeness principle of VM0051 \S9.3~\cite{vm0051}.
Third, the $72$-hour dry-phase counter in \eqref{eq:awd} is the operational
rule of Decision 4801 restated as a predicate on the replayed trajectory, so
compliance is evaluated rather than declared.

The system does not quantify CH$_4$ or N$_2$O, does not replace the emission
factors developed for Vietnamese cropping seasons, and does not replace
third-party verification. It reduces the cost of producing the evidence that
verification consumes.

\section{Discussion}
\label{sec:discussion}

\subsection{Relation to the command-gating twin}

The companion work~\cite{twingate2026} uses the same FAO-56 twin as a
\emph{pre-transmission} safety gate for irrigation commands over a lossy
Gilbert-Elliott channel~\cite{gilbert1960,elliott1963}, and proves a fidelity
threshold below which the gate is empty and silence is optimal. The present
paper uses the identical model as a \emph{post-hoc} auditor of what was claimed.
The two roles are complementary and share one code base: where the gate asks
``may I send this command'', the auditor asks ``could this command have
happened''. Deploying both on the same gateway yields the command record and its
attestation from one device, which is why the marginal cost of MRV in tier~0 is
close to zero.

\subsection{Why the filters are not merely a model check}

A reviewer may object that replaying a log through a model tests the model as
much as the log. Three replies. The forcing is public reanalysis that no
participant controls. The soil parameters are published averages for the delta
rather than fitted per farm~\cite{fao56}. And the guarantees are one-sided by
construction: Prop.~\ref{prop:noalarm} bounds the harm to honest farmers
independently of model error, because the tolerance $\varepsilon$ absorbs it.
Model error can therefore reduce detection power but cannot manufacture a false
accusation, which is the asymmetry that makes the layer deployable.

\subsection{Where remote sensing fits}

Satellite radar detects irrigation events but revisits every two to four days
and cannot resolve a $72$-hour dry phase in mid-period, which is precisely the
quantity that determines credit eligibility. Remote sensing is therefore used
here as a per-plot cross-check that lowers the frequency of field verification,
consistent with the digital MRV guidance of VM0051 Appendix 4~\cite{vm0051}, and
not as a substitute for the log. Existing AWD monitoring based on manual
observation tubes, low-cost wireless sensors or satellite
classification~\cite{delacruz2022,islam2026alose,lovell2019,mlrice2025} does not
quantify the reliability of the action record itself, which is the gap
identified by the survey underlying this work.

\subsection{Limitations and future work}

\begin{enumerate}
\item \emph{Simulation-based validation.} Fraud is injected into simulated
reference trajectories; a field pilot with real farmers and a real verifier is
needed to measure behavioural response, in particular whether farmers alter
behaviour once the filters are known.
\item \emph{Device prices are planning estimates.} The break-even condition
\eqref{eq:breakeven} is robust to them but the absolute margin is not; a
quotation must replace the estimates.
\item \emph{One shared pump is assumed.} The tier~1 closure of
Prop.~\ref{prop:omission} requires that the cooperative draws water through a
measurable point. Farms on tidal canals with individual inlets need one meter
per inlet, which raises $n_{\mathrm{dev}}$ and must be re-costed.
\item \emph{No emission quantification.} The layer produces activity data only.
Coupling it to a methane model and to locally calibrated emission factors is
separate work.
\item \emph{Salinity is a scenario.} The tidal-salinity constraint carried over
from the companion work has no field salinity measurements behind it and is
used only in sensitivity analysis.
\item \emph{Fixed soil parameters.} Using FAO-56 table averages for all three
stations ignores within-delta variability; a per-plot calibration step would
tighten $\varepsilon$ and raise detection power.
\end{enumerate}

The most valuable next step is a field pilot on one cooperative, because it
converts the two estimates that dominate this paper, device price and farmer
behaviour, into measurements.

\section{Conclusion}

We have shown that the evidence required to issue a rice carbon credit can be
produced by a farmer-written irrigation log attested by physics, at a marginal
monitoring cost near zero. Replaying the log through the FAO-56 water balance
under public ERA5 forcing exposes three violations that a compliant log cannot
produce, and we proved that no compliant log is ever accused, with a
closed-form bound of $61$\,mm on the largest claimable depth at the drainage
threshold. Omission is the one fraud class that physics cannot see; we proved
that a single flow meter at the cooperative pumping station closes it exactly,
and measured a gap of $29.0$\,mm against a $5$\,mm threshold. Across $240$
audited logs on three stations and two climate years the filters raised no
false alarm and detected every log of all three fraud classes. Expressing the
criterion in depletion fraction rather than absolute depth makes it invariant
to the $1.1$--$2.5$\,cm per year subsidence of the delta, which would otherwise
shift a $15$\,cm threshold by up to $83\%$ within a single crediting period.
The layer supplies activity data of the highest evidentiary rank under Decision
4801 and does not replace emission quantification or third-party verification.
Code, data and the invariant suite are public.

\appendix

\begin{table}[!t]
\centering
\caption{Experimental configuration. Every value is reproducible from the
released code and the public ERA5 forcing.}
\label{tab:setup}
\footnotesize
% Sinh TU DONG tu outputs/fraud_raw.csv boi experiments/make_tables_paper2.py
% KHONG sua tay: moi con so deu truy nguon duoc ve CSV.

\begin{tabular}{@{}p{2.6cm}p{5.9cm}@{}}
\toprule
Component & Value \\
\midrule
Weather & Hourly ERA5 via the Open-Meteo archive, no missing values \\
Stations & 3: Can Tho, Soc Trang, Ca Mau \\
Climate years & 2024, 2025 (2024 water-scarce, 2025 water-surplus) \\
Window & 1 March to 29 May (2\,160 h), the dry season of the Mekong Delta \\
Soil & Silty clay loam, FAO-56 Table 2: TAW $=70$, FC $=160$, WP $=90$\,mm \\
True irrigation schedule & AWD-compliant, irrigating at $d\ge 0.80$ with depths $\{20,40\}$\,mm \\
Log variants & 4 (one honest, three fraudulent) \\
Random seeds & 10 per cell \\
Logs audited in total & 	extbf{240} \\
Filters & $V_1,V_2,V_4$ physical (tier 0); pump-meter gap $G$ (tier 1) \\
\bottomrule
\end{tabular}

\end{table}

\begin{table}[!t]
\centering
\caption{Detection outcome per station and climate year. ``False alarm'' counts
honest logs wrongly flagged; the other three columns count fraud logs correctly
flagged out of $n=10$ per cell. Omission is detected at tier~1 by the pump-gap
test of \eqref{eq:gap}.}
\label{tab:detect}
\footnotesize
% Sinh TU DONG tu outputs/fraud_raw.csv boi experiments/make_tables_paper2.py
% KHONG sua tay: moi con so deu truy nguon duoc ve CSV.

\begin{tabular}{@{}p{1.5cm}p{0.75cm}p{0.35cm}p{1.25cm}p{1.25cm}p{1.2cm}p{1.2cm}@{}}
\toprule
Station & Year & $n$ & False alarm & Added events & False depth & Omission \\
\midrule
Ca Mau & 2024 & 10 & 0 & 10 & 10 & 10 \\
Ca Mau & 2025 & 10 & 0 & 10 & 10 & 10 \\
Can Tho & 2024 & 10 & 0 & 10 & 10 & 10 \\
Can Tho & 2025 & 10 & 0 & 10 & 10 & 10 \\
Soc Trang & 2024 & 10 & 0 & 10 & 10 & 10 \\
Soc Trang & 2025 & 10 & 0 & 10 & 10 & 10 \\
\midrule
\textbf{Total} & & 60 & 0 & 60 & 60 & 60 \\
\bottomrule
\end{tabular}

\end{table}

\begin{table}[!t]
\centering
\caption{Violation intensity per log type. $|L|$ is the number of claimed
events, $\sum u$ the claimed total depth, $|V_1|$ and $|V_4|$ the mean number of
saturation and storage-excess violations, $\bar d_{\max}$ the mean peak
depletion, and $G$ the pump gap of \eqref{eq:gap} in millimetres.}
\label{tab:intensity}
\footnotesize
% Sinh TU DONG tu outputs/fraud_raw.csv boi experiments/make_tables_paper2.py
% KHONG sua tay: moi con so deu truy nguon duoc ve CSV.

\begin{tabular}{@{}p{2.2cm}p{0.7cm}p{1.05cm}p{0.8cm}p{0.8cm}p{1.0cm}p{0.95cm}@{}}
\toprule
Log type & $|L|$ & $\sum u$ (mm) & $|V_1|$ & $|V_4|$ & $\bar d_{\max}$ & $G$ (mm) \\
\midrule
Honest log (AWD-compliant) & $4.3$ & $124.7$ & $0.00$ & $0.00$ & $0.802$ & $0.0$ \\
Fraud 1, added irrigation events & $12.3$ & $445.0$ & $5.87$ & $7.25$ & $0.802$ & $-320.3$ \\
Fraud 2, infeasible depth & $4.3$ & $307.0$ & $0.00$ & $1.57$ & $0.802$ & $-182.3$ \\
Fraud 3, omission to claim a dry phase & $3.3$ & $95.7$ & $0.00$ & $0.00$ & $0.896$ & $29.0$ \\
\bottomrule
\end{tabular}

\end{table}

\begin{table}[!t]
\centering
\caption{Derivation of the largest admissible claimed depth
(Prop.~\ref{prop:umax}) for the reference silty clay loam.}
\label{tab:prop1}
\footnotesize
% Sinh TU DONG tu outputs/fraud_raw.csv boi experiments/make_tables_paper2.py
% KHONG sua tay: moi con so deu truy nguon duoc ve CSV.

\begin{tabular}{@{}lll@{}}
\toprule
Quantity & Expression & Value \\
\midrule
$\mathrm{span}=FC-WP$ & $160-90$ & $70$\,mm \\
AWD drainage threshold & $d^{*}$ & $0.80$ \\
Deficit at $d^{*}$ & $d^{*}\cdot\mathrm{span}$ & $56.0$\,mm \\
Largest depth that never triggers $V_4$ & $d^{*}\cdot\mathrm{span}+\varepsilon$ & $61.0$\,mm \\
\bottomrule
\end{tabular}

\end{table}

\begin{table}[!t]
\centering
\caption{Tiered cost structure. Prices are planning estimates, not quotations;
the tier ordering and the break-even condition \eqref{eq:breakeven} do not
depend on them.}
\label{tab:cost}
\footnotesize
% Sinh TU DONG tu outputs/fraud_raw.csv boi experiments/make_tables_paper2.py
% KHONG sua tay: moi con so deu truy nguon duoc ve CSV.

\begin{tabular}{@{}p{0.7cm}p{2.6cm}p{1.2cm}p{3.6cm}@{}}
\toprule
Tier & Hardware & \$/ha/season & Additional evidence purchased \\
\midrule
0 & Irrigation log recorded on a mobile phone & $0$--$1$ & Tier-1 command record under Decision 4801 plus the physics filters \\
0+1 & One flow meter at the cooperative pumping station & $1$--$3$ & Detects omission fraud through the gap $G$ \\
0+1+2 & One water-level node shared by the whole cooperative & $6$--$11$ & Re-anchors the twin and satisfies the I7 calibration invariant \\
\bottomrule
\end{tabular}

\end{table}

\begin{table}[!t]
\centering
\caption{MRV cost per hectare per season for the four architectures. Monitoring
is the dominant term of MRV cost, so the architecture that avoids dedicated
sensing wins on labour rather than on hardware.}
\label{tab:baseline}
\footnotesize
% Sinh TU DONG tu outputs/fraud_raw.csv boi experiments/make_tables_paper2.py
% KHONG sua tay: moi con so deu truy nguon duoc ve CSV.

\begin{tabular}{@{}p{2.2cm}p{0.8cm}p{0.65cm}p{0.65cm}p{0.9cm}p{2.4cm}@{}}
\toprule
Architecture & Hardware & Labor & Audit & Total \$/ha & Note \\
\midrule
Manual observation tube plus photographs & $0.5$ & $8.0$ & $4$ & $12.5$ &current practice under Decision 4801 \\
Water-level sensor per plot & $11.3$ & $1.5$ & $4$ & $18.8$ &accurate, hard to scale to smallholders \\
Remote sensing plus field verification & $0$ & $2.0$ & $6$ & $8.0$ &revisit 2 to 4 days \\
\textbf{Physics-attested log (this work)} & $3.6$ & $0.3$ & $3$ &$\mathbf{8.4}$ & tier 0 alone is $\$0$--$1$ \\
\bottomrule
\end{tabular}

\end{table}

\begin{table}[!t]
\centering
\caption{Summary of headline results, each traced to its source.}
\label{tab:headline}
\footnotesize
% Sinh TU DONG tu outputs/fraud_raw.csv boi experiments/make_tables_paper2.py
% KHONG sua tay: moi con so deu truy nguon duoc ve CSV.

\begin{tabular}{@{}p{3.9cm}p{2.2cm}p{1.85cm}@{}}
\toprule
Result & Measured value & Source \\
\midrule
False alarms on honest logs & $0/60$ & Table~\ref{tab:detect} \\
Added irrigation events detected & 60/60 $=100\%$ & Table~\ref{tab:detect} \\
Infeasible depths detected & 60/60 $=100\%$ & Table~\ref{tab:detect} \\
Omission detected at tier 1 & 60/60 $=100\%$, $\bar G=29.0$\,mm & Table~\ref{tab:detect} \\
Largest depth that never raises an alarm & $u\le 61$\,mm & Prop.~\ref{prop:umax}, Table~\ref{tab:prop1} \\
Cost of tier 0 & $\sim\$0$--$1$/ha per season & Table~\ref{tab:cost} \\
Cost of the full stack (tiers 0+1+2) & $\sim\$6$--$11$/ha per season & Table~\ref{tab:cost} \\
Share of credit revenue & $5$--$15\%$ at $\$15$--$25$/t and $4$--$6$ credits/ha & Sec.~\ref{sec:cost} \\
Threshold drift from subsidence over a 5-year crediting period & $5.5$--$12.5$\,cm, that is $37$--$83\%$ of the 15\,cm AWD threshold & Prop.~\ref{prop:subsid} \\
\bottomrule
\end{tabular}

\end{table}

\begin{table*}[!t]
\centering
\caption{Mapping of the scheme onto Decision 4801/QD-BNNMT and VM0051 v1.1,
including what is explicitly not claimed.}
\label{tab:legal}
\footnotesize
% Sinh TU DONG tu outputs/fraud_raw.csv boi experiments/make_tables_paper2.py
% KHONG sua tay: moi con so deu truy nguon duoc ve CSV.

\begin{tabular}{@{}p{4.5cm}p{10.1cm}p{2.3cm}@{}}
\toprule
Requirement & How the system meets it & Level \\
\midrule
Decision 4801: the highest-ranked evidence is the record of actual irrigation and drainage actions &A signed, time-stamped command log at hourly resolution &Direct \\
Decision 4801: drainage counts only if the field is not re-irrigated for at least three days &A counter of dry phases of at least 72 h with $d\ge d_{\mathrm{dry}}$, evaluated on the replayed trajectory &Direct \\
Decision 4801, Appendix A12: uncertainty reported at a 95\% interval with $z=1.96$ &The same statistical machinery as the I7 calibration invariant of~\cite{twingate2026} &Direct \\
VM0051 \S 9.3: the principle of conservativeness when choosing parameters &Thresholds are biased toward raising an alarm ($\varepsilon=5$\,mm), so a compliant farmer may be queried but a fraudster is not passed &Aligned \\
VM0051 \S 5.2: a digital record of cultivation practices is acceptable &The mobile log is exactly this type of record &Direct \\
VM0051 Appendix 4: digital MRV and remote sensing are recommended &Sentinel-1 is used as a per-plot cross-check that lowers the frequency of field verification; it does not replace the log &Supporting \\
Not claimed &The system does not quantify CH$_4$ or N$_2$O (that requires the emission factors of IAE), does not replace third-party verification, and treats the salinity constraint as a scenario model only &Limitation \\
\bottomrule
\end{tabular}

\end{table*}

\section*{Data, Code and Declarations}
Code, the fraud-generation and replay pipeline, and all $240$ audited logs are
available at \url{https://github.com/mxuanvan02/TwinGate} under the MIT licence.
Weather forcing is public ERA5 reanalysis served by the Open-Meteo archive. This
research is funded by Hue University under Grant No. DHH2025-19-07. The authors
declare no competing interests. Ethics approval is not applicable: the study
uses public reanalysis data and simulation experiments, with no human
participants, animals or field interventions.

\bibliographystyle{IEEEtran}
\bibliography{refs}

\end{document}

